\documentclass[11pt,a4paper]{article}
\usepackage[margin=2.5cm]{geometry}
\usepackage[T1]{fontenc}
\usepackage{amsmath,amssymb,amsthm,mathtools}
\usepackage{enumitem}
\usepackage{microtype}
\usepackage{booktabs}
\usepackage{xcolor}
\usepackage{etoolbox}
\usepackage{hyperref}
\hypersetup{colorlinks=true,linkcolor=blue!50!black,citecolor=blue!50!black}

\newtheorem{theorem}{Theorem}[section]
\newtheorem{lemma}[theorem]{Lemma}
\newtheorem{proposition}[theorem]{Proposition}
\newtheorem{corollary}[theorem]{Corollary}
\theoremstyle{definition}
\newtheorem{definition}[theorem]{Definition}
\theoremstyle{remark}
\newtheorem{remark}[theorem]{Remark}

\theoremstyle{plain}
\newcommand{\restatedname}{}
\newtheorem*{restatedthm}{\restatedname}
\newenvironment{restate}[2]
  {\renewcommand{\restatedname}{#1}\begin{restatedthm}[#2]}
  {\end{restatedthm}}
\theoremstyle{definition}
\newtheorem*{restateddefn}{\restatedname}
\newenvironment{restatedef}[2]
  {\renewcommand{\restatedname}{#1}\begin{restateddefn}[#2]}
  {\end{restateddefn}}

%%%%%%%%%%%%%%%%%%%%%%%%%%%%%%%%%%%%%%%%%%%%%%%%%%%%%%%%%%%%%%%%%%%%%%%%%%%%%%
%%  notations.tex
%%  The 3/2 Phase Transition in Semi-Random 2-SAT
%%
%%  Symbols are organised by SCOPE, as variables are in a program:
%%
%%    Part I    general machinery (delimiters, operators, utilities)
%%    Part II   probability and expectation
%%    Part III  GLOBAL symbols: defined in Section 1, valid everywhere
%%    Part IV   MODULE symbols of Section 2 (graph-theoretic analysis)
%%    Part V    MODULE symbols of Section 3 (the main identity)
%%    Part VI   MODULE symbols of Section 4 and Appendices A, B
%%
%%  Rules.
%%   1. A global symbol is never reused for anything else.  Every symbol that
%%      appears in a statement of Section 1 is global.
%%   2. A module symbol is valid only in its own section (and appendices).  It
%%      may not reuse the letter of a global symbol.  Two different modules may
%%      reuse a letter, since their scopes do not overlap.
%%   3. A symbol that lives inside one proof (a local symbol) is written by
%%      hand in that proof and is not defined here.  It may not reuse the
%%      letter of a global symbol or of a module symbol of the same module.
%%   4. A subscript or a prime may specialise a symbol but not change its
%%      meaning: R_0, R_1 and R' are radii, p_0 is a retention probability.
%%   5. A symbol is defined only if it is used at least twice, or appears in
%%      a statement.  Everything else is written out in words.
%%
%%  checkscope.py checks rules 1--4 mechanically.  To rename a global or
%%  module symbol, edit the single line that defines it and recompile.
%%%%%%%%%%%%%%%%%%%%%%%%%%%%%%%%%%%%%%%%%%%%%%%%%%%%%%%%%%%%%%%%%%%%%%%%%%%%%%


%%%%%%%%%%%%%%%%%%%%%%%%%%%%%%%%%%%%%%%%%%%%%%%%%%%%%%%%%%%%%%%%%%%%%%%%%%%%%%
%% Part I.  General machinery
%%%%%%%%%%%%%%%%%%%%%%%%%%%%%%%%%%%%%%%%%%%%%%%%%%%%%%%%%%%%%%%%%%%%%%%%%%%%%%

% Delimiters that grow with their content in display style but stay small
% inline.  Everything below is built from these, so the paper never needs a
% hand-written \left...\right.
\newcommand{\AutoAdjust}[3]{%
  \mathchoice{\left#1#2\right#3}{#1#2#3}{#1#2#3}{#1#2#3}}
\newcommand{\InParentheses}[1]{\AutoAdjust{(}{#1}{)}}
\newcommand{\InBrackets}[1]{\AutoAdjust{[}{#1}{]}}
\newcommand{\InBraces}[1]{\AutoAdjust{\{}{#1}{\}}}

\newcommand{\bk}[1]{\InParentheses{#1}}
\newcommand{\Bk}[1]{\InBrackets{#1}}
\newcommand{\BK}[1]{\InBraces{#1}}
\newcommand{\abs}[1]{\left|{#1}\right|}
\newcommand{\norm}[1]{\left\lVert#1\right\rVert}
\newcommand{\ceil}[1]{\left\lceil{#1}\right\rceil}
\newcommand{\floor}[1]{\left\lfloor{#1}\right\rfloor}
\newcommand{\angbk}[1]{\left\langle{#1}\right\rangle}

% Sets and operators
\newcommand{\Z}{\mathbb{Z}}
\newcommand{\R}{\mathbb{R}}
\newcommand{\N}{\mathbb{N}}
\newcommand{\eps}{\varepsilon}
\newcommand{\dd}{\mathrm{d}}
\DeclareMathOperator*{\argmin}{arg\,min}
\DeclareMathOperator*{\argmax}{arg\,max}

\newcommand{\eqdef}{\overset{\mathrm{def}}{=\mathrel{\mkern-3mu}=}}
\newcommand{\defeq}{\eqdef}

% Utilities
\newcommand{\bnote}[1]{\textcolor{blue}{[#1]}}
\newcommand{\numberthis}{\addtocounter{equation}{1}\tag{\theequation}}
\newenvironment{proofof}[1]{\begin{proof}[Proof of #1]}{\end{proof}}


%%%%%%%%%%%%%%%%%%%%%%%%%%%%%%%%%%%%%%%%%%%%%%%%%%%%%%%%%%%%%%%%%%%%%%%%%%%%%%
%% Part II.  Probability and expectation
%%
%%  \E and \Prob are the bare symbols, for use as operators: \E\cyc{\Akept}.
%%  \Ex{...} and \Prx{...} carry their own delimiters and grow with content.
%%  Both take an optional subscript: \Ex[\Phi]{...}, \Prx[\mathrm{nb}]{...}.
%%
%%  To switch the whole paper between Pr(...) and Pr[...], or between an
%%  upright Pr and \mathbb P, edit only the four lines marked [STYLE].
%%%%%%%%%%%%%%%%%%%%%%%%%%%%%%%%%%%%%%%%%%%%%%%%%%%%%%%%%%%%%%%%%%%%%%%%%%%%%%

\newcommand{\E}{\mathbb{E}}                     % [STYLE] expectation symbol
\DeclareMathOperator{\Prob}{Pr}                 % [STYLE] probability symbol
% The optional argument is a subscript; \ifblank keeps an empty one from
% typesetting an empty subscript box.  The \! cancels the thin space TeX puts
% between the symbol and the delimiter, which after an operator such as Pr is
% wide enough to read as a gap.
\newcommand{\Ex}[2][]{\E\ifblank{#1}{}{_{#1}}\!\InBrackets{#2}}        % [STYLE] E[ ... ]
\newcommand{\Prx}[2][]{\Prob\ifblank{#1}{}{_{#1}}\!\InParentheses{#2}} % [STYLE] Pr( ... )
\newcommand{\Var}[2][]{\operatorname{\mathbf{Var}}\ifblank{#1}{}{_{#1}}\InBrackets{#2}}

% Indicator.  \1 is kept as a short alias for the bare symbol, so that the
% existing subscript form \1_{\{...\}} continues to work.
\newcommand{\1}{\mathbbm{1}}
\newcommand{\indic}[1]{\1\InBraces{#1}}

% Named events and distributions
\newcommand{\Bin}{\operatorname{Bin}}
\newcommand{\event}[1]{\mathcal{#1}}


%%%%%%%%%%%%%%%%%%%%%%%%%%%%%%%%%%%%%%%%%%%%%%%%%%%%%%%%%%%%%%%%%%%%%%%%%%%%%%
%% Part III.  GLOBAL symbols
%%
%%  Defined in Section 1 and listed in its notation table
%%  (Section~\ref{subsec:notation}).  Valid from the first page to the last.
%%%%%%%%%%%%%%%%%%%%%%%%%%%%%%%%%%%%%%%%%%%%%%%%%%%%%%%%%%%%%%%%%%%%%%%%%%%%%%

%% --- The input and its random formula -------------------------------------
\newcommand{\Gin}{G}                        % the input finite multigraph
\newcommand{\Vset}{V}                       % its vertex set
\newcommand{\Eset}{E}                       % its edge set
\newcommand{\nvert}{n}                      % |V|
\newcommand{\medge}{m}                      % |E|
\newcommand{\bvar}[1]{x_{#1}}               % Boolean variable at a vertex
\newcommand{\rform}[1]{\Phi(#1)}            % the random 2-CNF on a graph
\newcommand{\fml}{\varphi}                  % a fixed (non-random) 2-CNF formula
\DeclareMathOperator{\Sol}{Sol}             % solution set of a formula
\newcommand{\SAT}{\text{ SAT}}              % the event "... is satisfiable"
\newcommand{\indsat}[1][G]{\1_{_{\text{SAT}}}(\Phi(#1))}

%% --- Graph functionals ----------------------------------------------------
\newcommand{\ncomp}[1]{c(#1)}               % number of connected components
\newcommand{\cyc}[1]{\nu(#1)}               % cycle rank |E|-|V|+c
\newcommand{\exc}[1]{\xi(#1)}               % excess |E| - (3/2)|V|
\newcommand{\dg}[2]{\deg_{#1}(#2)}          % degree of a vertex
\newcommand{\dist}{\operatorname{dist}}     % graph distance

%% --- Balls ----------------------------------------------------------------
\newcommand{\ball}[2]{B_{#1}(#2)}           % radius-#1 ball around #2
\newcommand{\ballsz}[1]{M_{#1}}             % bound on the size of a radius-#1 ball

%% --- Solution graph -------------------------------------------------------
%  Sol(phi) is viewed as a graph (assignments at Hamming distance 1 are
%  joined), so c(.) has one meaning only: components of a graph.
\newcommand{\solcomp}[1]{\ncomp{\Sol(#1)}}  % number of solution clusters
\newcommand{\freeset}{S}                    % the set of free variables of a subcube
\newcommand{\fixasg}{\tau}                  % the values of the fixed variables
\newcommand{\cplx}[1]{X(#1)}                % cubical complex of subcubes inside Sol
\newcommand{\euler}[1]{\chi(#1)}            % Euler characteristic

%% --- The random edge-deletion experiment ----------------------------------
\newcommand{\Tdef}{T}                       % number of deferred edges
\newcommand{\Kdef}{K}                       % forest components holding a deferred edge
\newcommand{\Ffor}{F}                       % the forest built in the first stage

%% --- Percolation and long paths -------------------------------------------
\newcommand{\perc}[2]{#1_{#2}}              % H_p: keep each edge independently w.p. p
\newcommand{\keepp}{p}                      % retention probability
\newcommand{\plen}{\ell}                    % length of the paths we track
\newcommand{\plenz}{\ell_0}                 % the length finally selected
\newcommand{\Lpath}[3]{\mathcal{L}_{#1}(#2,#3)} % P(random root starts a length-#1 path)

%% --- Parameters -----------------------------------------------------------
\newcommand{\excpar}{\varepsilon}           % the excess parameter, xi(G) >= eps n
\newcommand{\Hyp}{\textup{(H)}}             % the standing hypothesis
% \loctree carries its own math delimiters and is used in TEXT mode, so that
% the words can be broken across lines.
\newcommand{\loctree}[2]{\((#1,#2)\)-locally tree-like}
\newcommand{\degb}{\Delta}                  % maximum-degree bound
\newcommand{\excd}{\delta}                  % excess density, xi(G) >= delta n
%  R is THE letter for a radius.  R and rho are the parameters of
%  Definition 1.1; R_0, R_1 are particular radii, R' is a radius inside a proof.
\newcommand{\rad}{R}                        % radius in the local-tree condition
\newcommand{\radz}{R_0}                     % threshold for R (Theorem 1.12)
\newcommand{\radp}{R_1}                     % radius of the packed balls (Lemmas 1.7, 1.8)
\newcommand{\nontree}{\rho}                 % tolerated fraction of non-tree balls
\newcommand{\nontreez}{\rho_0}              % threshold used for the bound on T
\newcommand{\nontreeo}{\rho_1}              % threshold of Theorem 1.12
\newcommand{\rate}{\gamma}                   % E[(3/4)^T] <= 2 e^{-gamma n}
\newcommand{\pack}{\eta}                    % fraction of vertices with a rich ball


%%%%%%%%%%%%%%%%%%%%%%%%%%%%%%%%%%%%%%%%%%%%%%%%%%%%%%%%%%%%%%%%%%%%%%%%%%%%%%
%% Part IV.  MODULE symbols of Section 2
%%
%%  None.  Every symbol of Section 2 is either global or local to one proof.
%%%%%%%%%%%%%%%%%%%%%%%%%%%%%%%%%%%%%%%%%%%%%%%%%%%%%%%%%%%%%%%%%%%%%%%%%%%%%%


%%%%%%%%%%%%%%%%%%%%%%%%%%%%%%%%%%%%%%%%%%%%%%%%%%%%%%%%%%%%%%%%%%%%%%%%%%%%%%
%% Part V.  MODULE symbols of Section 3
%%%%%%%%%%%%%%%%%%%%%%%%%%%%%%%%%%%%%%%%%%%%%%%%%%%%%%%%%%%%%%%%%%%%%%%%%%%%%%

\newcommand{\face}[2]{Q(#1,#2)}             % subcube (face) of the Boolean cube
\newcommand{\edgefac}[2]{w_{#1}(#2)}        % edge factor w_S(e): 3/4, 1/2 or 0


%%%%%%%%%%%%%%%%%%%%%%%%%%%%%%%%%%%%%%%%%%%%%%%%%%%%%%%%%%%%%%%%%%%%%%%%%%%%%%
%% Part VI.  MODULE symbols of Section 4 and Appendices A, B
%%
%%  The constants below are produced in exactly this order, and none of them
%%  depends on R or rho, which are chosen last.
%%%%%%%%%%%%%%%%%%%%%%%%%%%%%%%%%%%%%%%%%%%%%%%%%%%%%%%%%%%%%%%%%%%%%%%%%%%%%%

\newcommand{\Akept}{{\Ffor^{+}}}            % F plus the kept deferred edges (each edge w.p. 1/2)
\newcommand{\keeppz}{p_0}                   % the retention probability of Lemma 4.2
\newcommand{\plow}{\zeta}                   % lower bound on the long-path functional
\newcommand{\thr}{\kappa}                   % the error term of Lemma 4.3
\newcommand{\cycd}{\omega}                  % E nu(G_{1/2}) >= omega n

%% --- Appendix A: flows on paths -------------------------------------------
% \flowt carries its own math delimiters and is used in TEXT mode.
\newcommand{\flowt}[2]{\((#1,#2)\)-flow}
\newcommand{\brnum}{\lambda}                % branching rate demanded of good paths
\newcommand{\gfrac}{\psi}                   % fraction of good starting edges
\newcommand{\drift}{\mu}                    % mean drift of the log-branching walk
\newcommand{\ubd}{U}                        % upper bound on that walk's increments
\newcommand{\Pset}[1]{\mathcal{P}_{\plen}(#1)}  % candidate paths from a starting edge
\newcommand{\wt}[1]{w_{#1}}                 % the path weights
\newcommand{\pref}{\pi}                     % a path prefix

%% --- Appendix B: extinction in forests ------------------------------------
\newcommand{\sep}{\mathrm{Sep}}             % the separator set in a forest
\newcommand{\extp}[1]{h(#1)}                % forest extinction bound


\title{Reduction to difficult graphs via edge-disjoint barbells}
\date{}

\begin{document}
\maketitle

\noindent This note rewrites the graph-theoretic reduction (Section~1.4 and
Section~2 of the paper) along the plan in Section~\ref{sec:message}.
Section~\ref{sec:summary} summarizes the chapter, and
Section~\ref{sec:implementation} gives the statements and complete proofs. The
note is self-contained except for the final combination with Theorem~1.12 of
the paper (the analysis of difficult graphs), which is not changed.

\setcounter{section}{-1}
%======================================================================
\section{The plan (message of the advisor)}\label{sec:message}
%======================================================================

\begin{quote}\small
We can make the statement from cycles to barbells a purely graph-theoretic
statement: there is no need to assume that the cycles
are short. It should just be a statement that a connected subgraph with \(k\)
edge-disjoint cycles has a large number of barbells.

The next step is an easy analysis of how many edges the whole graph has, and
the conclusion that we must have enough short barbells, which leads to an
exponentially small probability of SAT.

Here is the plan for Step 1 of the proof:
\begin{enumerate}[label=\alph*)]
\item We eliminate all vertices of degree \(0\) and \(1\), (i) improving the
  excess, (ii) not changing the probability of SAT, and (iii) keeping the
  number of edges in \(G\) large. We then remove all connected components with
  excess \(\le0\). Again, the number of edges in the graph must remain large.
\item We reduce the largest degree \(\Delta(G)\) to a constant \(\Delta_0\)
  (which depends on the constant in the linear excess): we (i) may only
  increase the probability of SAT, (ii) add vertices and edges, and
  (iii) decrease the excess, but still keep it positive and linear in \(n\).
\item We state that graphs with many short edge-disjoint cycles are easy to
  analyze, and show that \(\Pr[\mathrm{SAT}]<e^{-\Omega(n)}\). We first
  prove/state a lemma. Lemma (graph theory): if a connected graph has \(k\)
  edge-disjoint cycles, then it has a large number of barbells with endpoints
  at those cycles.

  Then we state the corollary: if a graph has many (linear in \(n\))
  edge-disjoint short cycles and positive excess in every connected
  component, then \(\Pr[\mathrm{SAT}]<e^{-\Omega(n)}\). The proof of this
  corollary is easy: we look at all connected components and consider two
  cases: (i) only one short cycle in the component; (ii) \(\ge2\) short cycles
  in the component. In case (ii) we consider a minimal spanning subgraph
  connecting all those edge-disjoint cycles in the component and apply the
  graph-theoretic lemma. In case (i) we just find one barbell using only
  vertices of the connected component, which are disjoint from all remaining
  vertices. We then do a simple counting argument for the total number of
  edges in the graph and conclude that many of the barbells cannot have very
  long paths (otherwise there would be too many edges).
\item Finally, we can say that the graph can have only a sublinear number of
  loops (they form natural edge-disjoint short cycles) and a sublinear number
  of vertices with non-tree-like neighborhoods. Here the only claim to be
  stated and proved is that in bounded-degree graphs (\(\Delta(G)<\Delta_0\)),
  if there are \(\rho n\) vertices with non-tree-like neighborhoods of constant
  radius \(R\), then there is a constant \(C(\Delta_0,R)>0\) such that \(G\)
  has \(C\rho n\) edge-disjoint cycles. The proof of this claim is
  straightforward: consider all vertices with non-tree-like neighborhoods, and
  notice that each such cycle may intersect at most a constant number of other
  such cycles.
\end{enumerate}
\end{quote}

%======================================================================
\section{Summary of the chapter}\label{sec:summary}
%======================================================================

This section states all results of the chapter in the form in which they are
proved in Section~\ref{sec:implementation}; each statement is repeated there
word for word, with its proof. The notation is that of the paper and is
recalled in Section~\ref{subsec:setting}. In particular the degree bound
\(\Delta_0\) of the plan is the constant \(\degb\) of the paper, \(c(\cdot)\)
always counts components, and \Hyp{} is the hypothesis: no self-loops, all
degrees in \([2,\degb]\), and \(\exc{H}\ge\excd|V(H)|\). A graph is
\emph{difficult} if it satisfies \Hyp{} and is \loctree{\rad}{\nontree}.

\subsection{The main result}

\begin{restate}{Theorem~\ref{thm:reduction}}{Reduction to difficult graphs}
Fix \(\excpar>0\), and let \(\degb,\excd\) be as in
Proposition~\ref{prop:cleaning}. For every integer \(\rad\ge1\) and every
\(\nontree\in(0,1]\) there is \(\beta=\beta(\excpar,\rad,\nontree)>0\) such that for
every multigraph \(\Gin\) with \(\nvert\) vertices and \(\exc{\Gin}\ge\excpar \nvert\), one
of the following holds:
\begin{enumerate}[label=(\alph*)]
\item \(\Prx{\rform{\Gin}\SAT}\le e^{-\beta \nvert}\);
\item there is a multigraph \(H\) that satisfies \Hyp{} with
  \(\degb,\excd\), is \loctree{\rad}{\nontree}, has
  \(|V(H)|\ge\beta \nvert\), and satisfies
  \(\Prx{\rform{\Gin}\SAT}\le\Prx{\rform{H}\SAT}\).
\end{enumerate}
\end{restate}

\noindent Combined with the paper, case~(b) also gives
\(\Prx{\rform{\Gin}\SAT}\le e^{-\Omega(\nvert)}\). The paper's Theorem~1.12 gives
\(\radz(\degb,\excd)\) and \(\nontreeo(\degb,\excd)\); apply the theorem above
with \(\rad=\radz\) and \(\nontree=\nontreeo\). Then \(H\) is difficult. Since
\(\solcomp{\rform{H}}\ge1\) whenever \(\rform{H}\) is satisfiable,
the main identity (Theorem~1.10) together with Theorem~1.12 gives
\(\Prx{\rform{H}\SAT}\le\E\,\solcomp{\rform{H}}\le
2e^{-\rate|V(H)|/2}\), with \(|V(H)|\ge\beta \nvert\). There is no circularity,
because \(\radz\) and \(\nontreeo\) depend only on \(\degb\) and \(\excd\).

\subsection{Step 1: cleaning (plan (a) and (b))}

\begin{restate}{Lemma~\ref{lem:leaf}}{Deleting low-degree vertices}
Deleting an isolated vertex, or a vertex of degree \(1\) together with its
edge, does not change satisfiability and increases the excess by \(3/2\),
respectively \(1/2\).
\end{restate}

\begin{restate}{Lemma~\ref{lem:splitting}}{Splitting vertices}
Fix an integer \(\degb\ge4\). Every finite multigraph \(\Gin\) with \(\medge\)
edges can be changed into a finite multigraph \(H\) with maximum degree at
most \(\degb\) and with
\[
  N:=|V(H)|-|V(\Gin)|=|E(H)|-|E(\Gin)|\le\frac{2\medge}{\degb-2},
\]
in which no vertex has degree \(0\) or \(1\) unless it had this degree in
\(\Gin\). Moreover, there is a coupling of the random signs under which
\(\{\rform{\Gin}\SAT\}\subseteq\{\rform{H}\SAT\}\), and under which the signs on
\(H\) are independent and uniform, so that they have the law of \(\rform{H}\).
\end{restate}

\begin{restate}{Proposition~\ref{prop:cleaning}}{Cleaning}
Fix \(\excpar>0\), and let
\[
  \degb:=2+\Bigl\lceil\max\Bigl\{4,\frac6\excpar\Bigr\}\Bigr\rceil,
  \qquad
  \excd:=\frac{\excpar}{2(4+\excpar)} .
\]
For every multigraph \(\Gin\) with \(\nvert\) vertices and \(\exc{\Gin}\ge\excpar \nvert\)
there is a multigraph \(H_1\) such that
\begin{enumerate}[label=(\roman*)]
\item \(2\le\dg{H_1}{v}\le\degb\) for every vertex \(v\);
\item every component of \(H_1\) has positive excess;
\item \(\exc{H_1}\ge2\excd|V(H_1)|\);
\item \(|V(H_1)|\ge(\excpar/\degb)\,\nvert\);
\item \(\Prx{\rform{\Gin}\SAT}\le\Prx{\rform{H_1}\SAT}\).
\end{enumerate}
\end{restate}

\noindent The graph \(H_1\) is obtained by splitting first, then deleting
vertices of degree \(0,1\) (Lemma~\ref{lem:leaf}), and then deleting components of
nonpositive excess.

\subsection{Step 2: many edge-disjoint cycles (plan (c))}

\begin{restatedef}{Definition~\ref{def:barbell}}{Barbell}
A \emph{barbell} is a triple \(Y=(D_1,D_2,P)\), where \(D_1,D_2\) are
edge-disjoint cycles, its \emph{ends}, and \(P\) is a path from a vertex of
\(D_1\) to a vertex of \(D_2\) that uses no edge of \(D_1\) or \(D_2\). Its edge
set is \(E(Y)=E(D_1)\cup E(D_2)\cup E(P)\). If \(D_1\) and \(D_2\) share a
vertex, \(P\) may be this single vertex.
\end{restatedef}

\noindent Here a cycle may be a self-loop or a pair of parallel edges.

\begin{restate}{Lemma~\ref{lem:forcing}}{Forcing cycles}
Let \(C\) be a cycle and \(u\) a vertex of \(C\). There is a sign pattern on
\(E(C)\) for which every satisfying assignment of the clauses on \(E(C)\) has
\(\bvar{u}=1\), and one for which every satisfying assignment has \(\bvar{u}=0\).
\end{restate}

\begin{restate}{Lemma~\ref{lem:barbell-unsat}}{Barbells give contradictions}
Every barbell supports an unsatisfiable 2-CNF. Hence, if a barbell \(Y\) is a
subgraph of \(\Gin\), then with probability at least \(4^{-|E(Y)|}\) the signs on
\(E(Y)\) make \(\rform{\Gin}\) unsatisfiable.
\end{restate}

\noindent This is the letter Q: one cycle forces the value \(1\) at the start
of the path (Lemma~\ref{lem:forcing}), the clauses on the path carry it to the other
cycle, and that cycle forces the value \(0\) there.

\begin{restate}{Lemma~\ref{lem:walk}}{From walks to barbells}
Let \(D_1,D_2\) be edge-disjoint cycles, and let \(W\) be a walk from a vertex
of \(D_1\) to a vertex of \(D_2\) that uses no edge of \(D_1\) or \(D_2\). Then
there is a path \(P\) with \(E(P)\subseteq E(W)\) such that \((D_1,D_2,P)\) is
a barbell.
\end{restate}

\begin{restate}{Lemma~\ref{lem:excess-cycles}}{Positive excess gives two cycles}
Every finite connected multigraph \(\Gin\) with \(\exc{\Gin}>0\) contains two
edge-disjoint cycles.
\end{restate}

\begin{restate}{Corollary~\ref{cor:excess-barbell}}{Positive excess gives a barbell}
Every finite connected multigraph \(\Gin\) with \(\exc{\Gin}>0\) contains a barbell.
\end{restate}

\noindent The corollary follows by joining the two cycles of Lemma~\ref{lem:excess-cycles} by a
shortest path between them. It handles the components that
contain only one of the given cycles, in particular graphs with many
components, each with a single short cycle. The hypothesis is positive
excess, not positive cycle rank: a theta graph has \(\nu=2\) and no two
edge-disjoint cycles.

\begin{restate}{Theorem~\ref{thm:barbell}}{Barbell lemma}
Let \(\Gin\) be a connected multigraph that contains pairwise edge-disjoint
cycles \(C_1,\dots,C_k\). Then \(\Gin\) contains at least
\(\lceil(k-1)/3\rceil\) pairwise edge-disjoint barbells, all of whose ends are
among \(C_1,\dots,C_k\).
\end{restate}

\noindent This is the purely graph-theoretic statement asked for in the plan:
no length condition on the cycles. The proof is a bottom-up pairing on a
spanning tree in which each cycle is contracted to a node. Every subtree
passes at most one unmatched cycle upward, and ordinary nodes pair the
arriving cycles. A cycle node with zero arrivals passes its own cycle up;
with one arrival it pairs the arrival with its own cycle; with at least two
arrivals it pairs them through edge-disjoint arcs of its cycle and discards
its own cycle. At the end every cycle is an end of one of the \(b\)
barbells (at most \(2b\) cycles), the one cycle left unmatched at the root (at
most \(1\)), or discarded. Every discarded cycle is charged to a barbell
created at the same node, so at most \(b\) cycles are discarded. Hence
\(k\le2b+1+b\). The bound is sharp: join a vertex by an edge to each of \(t\)
gadgets, each gadget being a triangle (entered at one vertex) with one more
triangle hanging at each of its two other vertices. This gives \(k=3t\)
edge-disjoint triangles. At most two triangles of each gadget can be ends:
a barbell with the middle triangle as an end uses all of its edges, and an
outer triangle that is not its partner can reach any other cycle only through
an edge of the middle triangle. So at most \(2t\) triangles are ends, and
there are at most \(t\) edge-disjoint barbells.

\begin{restate}{Corollary~\ref{cor:many-cycles}}{Many edge-disjoint cycles}
Let \(\Gin\) be a multigraph with \(\nvert\) vertices and maximum degree at most
\(\degb\) in which every component has positive excess. If \(\Gin\) contains at
least \(\alpha \nvert\) pairwise edge-disjoint cycles, where \(\alpha>0\), then
\[
  \Prx{\rform{\Gin}\SAT}\le\exp\Bigl(-\frac\alpha8\,4^{-4\degb/\alpha}\,\nvert\Bigr).
\]
\end{restate}

\noindent Proof idea: let a component contain \(k_j\) of the cycles. If
\(k_j\ge2\), Theorem~\ref{thm:barbell} gives
\(\lceil(k_j-1)/3\rceil\ge k_j/4\) edge-disjoint barbells; the inequality is
\(1\ge k_j/4\) for \(k_j=2,3\), and \(4(k_j-1)\ge3k_j\) for \(k_j\ge4\). If
\(k_j=1\), Corollary~\ref{cor:excess-barbell} gives one barbell, and
\(1\ge k_j/4\). Summing over components gives at least \(\alpha \nvert/4\)
edge-disjoint barbells. The constant \(1/4\) is the best uniform one, since
\(k_j=4\) only guarantees one barbell. Since \(|E(\Gin)|\le\degb \nvert/2\), at least \(\alpha \nvert/8\) of these
barbells have at most \(4\degb/\alpha\) edges, and they fail independently.

\subsection{Step 3: few cycles (plan (d))}

\begin{restate}{Lemma~\ref{lem:nontree-cycles}}{Non-tree balls give disjoint cycles}
Let \(\Gin\) have maximum degree at most \(\degb\), let \(\rad\ge1\), and suppose
that \(N\) vertices \(v\) have a ball \(\ball{\rad}{v}\) that is not a tree. Then \(\Gin\)
contains at least \(N/\ballsz{2\rad}\) pairwise vertex-disjoint cycles, each of length
at most \(2\rad+1\).
\end{restate}

\noindent This is the claim of plan (d), with \(C(\Delta_0,\rad)=1/\ballsz{2\rad}\).

\begin{restate}{Lemma~\ref{lem:remove-loops}}{Removing self-loops}
Fix \(\degb\), \(\excd>0\), \(\rad\ge1\) and \(\nontree\in(0,1]\). Let \(\Gin\) have
\(\nvert\) vertices, \(2\le\dg{\Gin}{v}\le\degb\) for all \(v\), and
\(\exc{\Gin}\ge2\excd \nvert\); suppose that \(\Gin\) has at most \(\excd \nvert\)
self-loops and at most \((2\excd\nontree/\degb)\nvert\) vertices with a non-tree
ball \(\ball{\rad}{v}\). Delete all self-loops, then repeatedly delete vertices of
degree \(0\) or \(1\), and call the result \(H\). Then \(H\) satisfies \Hyp{}
with \(\degb,\excd\), it is \loctree{\rad}{\nontree},
\(|V(H)|\ge(2\excd/\degb)\nvert\), and
\(\Prx{\rform{\Gin}\SAT}\le\Prx{\rform{H}\SAT}\).
\end{restate}

\subsection{How the statements give the main result}

Theorem~\ref{thm:reduction} has two possible conclusions: (a) the
satisfiability probability is exponentially small, or (b) it is bounded by
that of a difficult graph \(H\) with \(|V(H)|\ge\beta \nvert\). Which one we get
depends on the number of edge-disjoint cycles of the cleaned graph.

Let \(H_1\) be the graph of Proposition~\ref{prop:cleaning},
\(n_1=|V(H_1)|\ge(\excpar/\degb)\nvert\), and
\[
  \alpha:=\frac{2\excd\nontree}{\degb \ballsz{2\rad}}\le\excd .
\]
\begin{itemize}
\item \emph{\(H_1\) has at least \(\alpha n_1\) edge-disjoint cycles.} Then
  Corollary~\ref{cor:many-cycles} gives
  \(\Prx{\rform{H_1}\SAT}\le e^{-\Omega(n_1)}\). Since
  \(\Prx{\rform{\Gin}\SAT}\le\Prx{\rform{H_1}\SAT}\) and \(n_1\ge(\excpar/\degb)\nvert\),
  this is conclusion~(a).
\item \emph{\(H_1\) has fewer than \(\alpha n_1\) edge-disjoint cycles.}
  Self-loops are edge-disjoint cycles, so \(H_1\) has fewer than
  \(\alpha n_1\le\excd n_1\) self-loops. By Lemma~\ref{lem:nontree-cycles} it
  has fewer than \(\alpha \ballsz{2\rad}n_1=(2\excd\nontree/\degb)n_1\) vertices with a
  non-tree ball. So Lemma~\ref{lem:remove-loops} applies to \(H_1\), and the
  resulting \(H\) is difficult, with
  \(|V(H)|\ge(2\excd/\degb)n_1\ge(2\excd\excpar/\degb^2)\nvert\). This is
  conclusion~(b).
\end{itemize}
The single threshold \(\alpha\) is chosen for two reasons. First,
\(\alpha\le\excd\), so the fewer than \(\excd n_1\) self-loops cost at most
half of the excess \(\exc{H_1}\ge2\excd n_1\), and \(\exc{H}\ge\excd n_1\)
survives. Second, \(\alpha \ballsz{2\rad}=2\excd\nontree/\degb\): the graph \(H\) may
have fewer vertices than \(H_1\), but \(|V(H)|\ge(2\excd/\degb)n_1\), so the
fewer than \((2\excd\nontree/\degb)n_1\) non-tree balls are still at most
\(\nontree|V(H)|\). The constants are
\[
  \beta=\min\Bigl\{\frac{\excpar}{\degb}\cdot\frac\alpha8\,
  4^{-4\degb/\alpha},\ \frac{2\excd\excpar}{\degb^2}\Bigr\}.
\]

\subsection{Changes with respect to the plan}
\begin{enumerate}[label=(\arabic*)]
\item \emph{Order of steps a) and b) of the plan.} Splitting a vertex lowers the excess of its
  component, so a component of small positive excess may become nonpositive.
  We therefore split first and clean afterwards. Splitting creates no vertex
  of degree \(0\) or \(1\), so the cleaning is done only once.
\item \emph{No dense case.} The paper uses the first-moment bound for
  \(\medge\ge2.41\nvert\) only to control the number of vertices added by the splitting.
  A direct computation (Proposition~\ref{prop:cleaning}) shows that the
  splitting works for every density, so the first-moment bound is no longer
  needed. The splitting adds at most \(2\medge/(\degb-2)\) vertices and lowers the
  excess by at most \(\medge/(\degb-2)\), where \(\medge=\frac32\nvert+\exc{\Gin}\). Both the
  \(\nvert\)-part and the \(\exc{\Gin}\)-part of this loss are small fractions of
  \(\exc{\Gin}\ge\excpar \nvert\) once \(\degb\ge2+6/\excpar\), however large
  \(\medge\) is. What has to stay linear is the excess \emph{relative to the new
  number of vertices}, because \Hyp{} asks for \(\exc{H}\ge\excd|V(H)|\).
\item \emph{``Short'' is not needed in Step 2.} Corollary~\ref{cor:many-cycles}
  assumes only linearly many edge-disjoint cycles. Bounded size of the barbells
  follows from the edge count \(|E|\le\degb|V|/2\).
\item \emph{Case (i) of step c) needs a lemma.} A component with only one of
  the chosen cycles still contains a barbell, but this uses its positive excess:
  Lemma~\ref{lem:excess-cycles} shows that a connected graph of positive
  excess contains two edge-disjoint cycles. This is where graphs with many
  components, each with a single short cycle, are handled. The minimal
  spanning subgraph in case (ii) is not needed, since the barbell lemma holds
  for every connected graph.
\item \emph{``Sublinear'' becomes quantitative.} Step 3 needs at most
  \(\alpha n_1\) self-loops and non-tree balls, with the explicit constant
  \(\alpha\) above; a genuine \(o(\nvert)\) bound is neither available nor needed.
\item \emph{After deleting the self-loops} one must delete the new vertices of
  degree \(0,1\), because \Hyp{} asks for minimum degree \(2\). Components of
  nonpositive excess need not be deleted again: \Hyp{} does not ask for
  positive excess per component.
\end{enumerate}

\subsection{What disappears from the paper}
The following parts of the paper
are no longer needed:
\begin{itemize}
\item the first-moment bound (Lemma~1.3) as a case of the proof; it can
  remain as motivation;
\item the cycle-rank criterion for unsatisfiable support (Lemma~1.6, proved
  by Lemmas~2.3 and~2.4), and with it the reference to Karve--Hirani;
\item the separate treatment of self-loops (Lemma~1.7 and Lemma~2.5);
\item the local-tree counting lemma (Lemma~1.8), with its radii \(s\),
  \(\radp=2s+\rad\) and three cases;
\item the packing of disjoint balls (Corollary~1.9).
\end{itemize}
Lemmas~1.4 and~1.5 (bounded degree and positive-excess components) are
replaced by Proposition~\ref{prop:cleaning}. The removed parts are replaced
by one graph-theoretic theorem (the barbell lemma) and one counting argument
(Corollary~\ref{cor:many-cycles}).

%======================================================================
\section{Implementation}\label{sec:implementation}
%======================================================================

\subsection{Setting}\label{subsec:setting}

A \emph{multigraph} may have parallel edges and self-loops. Each vertex \(v\)
carries a Boolean variable \(\bvar{v}\), and each edge \(uv\) carries one of the
four clauses \(\bvar{u}\vee \bvar{v}\), \(\bvar{u}\vee\neg \bvar{v}\), \(\neg \bvar{u}\vee \bvar{v}\),
\(\neg \bvar{u}\vee\neg \bvar{v}\): each end of the edge gets a sign, which decides
whether its variable appears negated. The two signs of an edge are independent
and uniform, and different edges are independent. A self-loop at \(v\) carries
a clause on \(\bvar{v}\) alone, again with two independent uniform signs. The
resulting random 2-CNF is \(\rform{\Gin}\). Deleting an
edge deletes its clause, which can only make the formula easier to satisfy.

For a multigraph \(H\) we write \(\exc{H}=|E(H)|-\frac32|V(H)|\) for its
\emph{excess}; it is additive over components. In the degree \(\dg{H}{v}\) a
self-loop at \(v\) adds \(2\). The ball \(\ball{\rad}{v}\) is the subgraph induced by
the vertices \(u\) with \(\dist(u,v)\le \rad\); it is connected. If the maximum
degree is at most \(\degb\), then \(|V(\ball{\rad}{v})|\le \ballsz{\rad}:=1+\degb\sum_{j=0}^{\rad-1}
(\degb-1)^j\). A graph is \emph{\loctree{\rad}{\nontree}} if at most
\(\nontree|V|\) vertices \(v\) have a ball \(\ball{\rad}{v}\) that is not a tree.

A \emph{cycle} of length \(L\ge1\) is a self-loop (\(L=1\)), two parallel edges
(\(L=2\)), or, for \(L\ge3\), a closed path through \(L\) distinct vertices. A
\emph{path} may consist of a single vertex. A multigraph \(H\) \emph{supports an
unsatisfiable 2-CNF} if for some sign pattern on \(E(H)\) the clauses on
\(E(H)\) are unsatisfiable.

The goal of the reduction is the following hypothesis, for constants
\(\degb\) and \(\excd>0\):

\medskip
\noindent\Hyp{} \(H\) has no self-loops, \(2\le\dg{H}{v}\le\degb\) for every
vertex, and \(\exc{H}\ge\excd|V(H)|\).
\medskip

\noindent A graph that satisfies \Hyp{} and is \loctree{\rad}{\nontree} is
\emph{difficult}.

%----------------------------------------------------------------------
\subsection{Step 1: bounded degree and positive excess}
%----------------------------------------------------------------------

\begin{lemma}[Deleting low-degree vertices]\label{lem:leaf}
Deleting an isolated vertex, or a vertex of degree \(1\) together with its
edge, does not change satisfiability and increases the excess by \(3/2\),
respectively \(1/2\).
\end{lemma}

\begin{proof}
An isolated variable does not occur in the formula. A variable that occurs in
one clause only can always be chosen to satisfy it. The changes of the excess
are \(\frac32\) and \(-1+\frac32=\frac12\).
\end{proof}

\begin{lemma}[Splitting vertices]\label{lem:splitting}
Fix an integer \(\degb\ge4\). Every finite multigraph \(\Gin\) with \(\medge\)
edges can be changed into a finite multigraph \(H\) with maximum degree at
most \(\degb\) and with
\[
  N:=|V(H)|-|V(\Gin)|=|E(H)|-|E(\Gin)|\le\frac{2\medge}{\degb-2},
\]
in which no vertex has degree \(0\) or \(1\) unless it had this degree in
\(\Gin\). Moreover, there is a coupling of the random signs under which
\(\{\rform{\Gin}\SAT\}\subseteq\{\rform{H}\SAT\}\), and under which the signs on
\(H\) are independent and uniform, so that they have the law of \(\rform{H}\).
\end{lemma}

\begin{proof}
Let \(b:=\degb-2\ge2\). For every vertex \(v\) with \(\dg{\Gin}{v}>b\), let
\(k_v=\lceil\dg{\Gin}{v}/b\rceil\), replace \(v\) by \(k_v\) copies, split the
\(\dg{\Gin}{v}\) edge ends at \(v\) into \(k_v\) nonempty groups of size at most
\(b\), give one group to each copy, and join the copies by a path. The number
of added vertices equals the number of added path edges, namely
\(N=\sum_v(k_v-1)\le\sum_v\dg{\Gin}{v}/b=2\medge/(\degb-2)\). Every copy has at most
\(b\) original edge ends and at most two path edges, so its degree is at most
\(\degb\), and since every group is nonempty and \(k_v\ge2\), every copy has
degree at least \(2\).

For the coupling, fix a split vertex with copies \(1,\dots,k_v\) along the
path. First put on each path edge \(ij\) both clauses
\((\neg \bvar{i}\vee \bvar{j})\wedge(\bvar{i}\vee\neg \bvar{j})\); they force all copies to be
equal, so the formula is satisfiable exactly when \(\rform{\Gin}\) is. Next choose
independent fair bits \(g_i\), one per copy, and replace \(\bvar{i}\) by
\(\bvar{i}\oplus g_i\) in every clause; this does not change satisfiability.
Finally keep on every path edge only one of its two clauses, chosen
independently and uniformly. The result is a subformula, so it is satisfiable
whenever \(\rform{\Gin}\) is.

It remains to show that the resulting signs are independent and uniform.
Consider one path with copies \(1,\dots,k_v\). Its randomness consists of the
bits \(g_1,\dots,g_{k_v}\) and the \(k_v-1\) choices of kept clauses, so
\(2^{k_v}\cdot2^{k_v-1}\) equally likely outcomes, and it produces the signs
on the \(k_v-1\) path edges, \(4^{k_v-1}\) possible patterns. After the xor,
the sign at end \(i\) of the edge \(\{i,i+1\}\) is the sign of \(\bvar{i}\) in the
kept clause flipped by \(g_i\), and likewise at end \(i+1\). So for fixed
\(g_i\), the target sign at end \(i\) determines which clause is kept, and
then the target sign at end \(i+1\) determines \(g_{i+1}\). Hence, once
\(g_1\) is fixed, every target pattern arises from exactly one outcome, so
every pattern arises from exactly two of the \(2\cdot4^{k_v-1}\) outcomes, and
the signs on the path edges are uniform. Different paths use independent
randomness. Finally, every sign at an end of an original edge is the original
sign flipped by the bit of its copy. The original signs are uniform and
independent of all bits and choices, so the flipped signs are again uniform
and independent of them, and hence of the signs on the path edges.
\end{proof}

\begin{proposition}[Cleaning]\label{prop:cleaning}
Fix \(\excpar>0\), and let
\[
  \degb:=2+\Bigl\lceil\max\Bigl\{4,\frac6\excpar\Bigr\}\Bigr\rceil,
  \qquad
  \excd:=\frac{\excpar}{2(4+\excpar)} .
\]
For every multigraph \(\Gin\) with \(\nvert\) vertices and \(\exc{\Gin}\ge\excpar \nvert\)
there is a multigraph \(H_1\) such that
\begin{enumerate}[label=(\roman*)]
\item \(2\le\dg{H_1}{v}\le\degb\) for every vertex \(v\);
\item every component of \(H_1\) has positive excess;
\item \(\exc{H_1}\ge2\excd|V(H_1)|\);
\item \(|V(H_1)|\ge(\excpar/\degb)\,\nvert\);
\item \(\Prx{\rform{\Gin}\SAT}\le\Prx{\rform{H_1}\SAT}\).
\end{enumerate}
\end{proposition}

\begin{proof}
Write \(\xi=\exc{\Gin}\) and \(\medge=|E(\Gin)|=\frac32\nvert+\xi\), and let
\(\theta:=1/(\degb-2)\), so that \(\theta\le\frac14\) and
\(\theta\le\excpar/6\).

\emph{Splitting.} Lemma~\ref{lem:splitting} gives \(H'\) with maximum degree
at most \(\degb\) and \(N\le2\theta \medge\) new vertices and edges. Each new
vertex with its new edge lowers the excess by \(\frac12\), so
\[
  \exc{H'}=\xi-\frac N2\ge\xi-\theta\Bigl(\frac32\nvert+\xi\Bigr)
  \ge\frac34\xi-\frac\excpar4\nvert\ge\frac\xi2,
\]
using \(\theta\le\frac14\), \(\frac32\theta\le\frac\excpar4\) and
\(\excpar \nvert\le\xi\). Also
\(|V(H')|=\nvert+N\le \nvert+2\theta(\frac32\nvert+\xi)\le2\nvert+\frac\xi2\), hence
\[
  \frac{\exc{H'}}{|V(H')|}\ge\frac{\xi/2}{2\nvert+\xi/2}=\frac{\xi}{4\nvert+\xi}
  \ge\frac{\excpar}{4+\excpar}=2\excd ,
\]
since the middle expression increases with \(\xi\ge\excpar \nvert\).

\emph{Cleaning.} Repeatedly delete vertices of degree \(0\) or \(1\), and then
delete every component of nonpositive excess; call the result \(H_1\). By
Lemma~\ref{lem:leaf} the first deletions do not lower the excess, and since the
excess is additive over components, neither do the second. The number of
vertices does not grow, so
\(\exc{H_1}\ge\exc{H'}\ge2\excd|V(H')|\ge2\excd|V(H_1)|\). Deletions do not
raise degrees, the first deletions stop only when all degrees are at least
\(2\), and deleting whole components does not change the remaining degrees;
this gives (i), and (ii) holds by construction. Since
\(\exc{H_1}\le|E(H_1)|\le\frac{\degb}2|V(H_1)|\) and
\(\exc{H_1}\ge\exc{H'}\ge\xi/2\ge\excpar \nvert/2\), we get (iv). Finally, (v)
follows from the coupling of Lemma~\ref{lem:splitting}, from
Lemma~\ref{lem:leaf}, and from the fact that deleting components deletes
clauses.
\end{proof}

%----------------------------------------------------------------------
\subsection{Local contradictions: the letter Q}
%----------------------------------------------------------------------

\begin{lemma}[Forcing cycles]\label{lem:forcing}
Let \(C\) be a cycle and \(u\) a vertex of \(C\). There is a sign pattern on
\(E(C)\) for which every satisfying assignment of the clauses on \(E(C)\) has
\(\bvar{u}=1\), and one for which every satisfying assignment has \(\bvar{u}=0\).
\end{lemma}

\begin{proof}
Write the cycle as \(u=u_0,u_1,\dots,u_L=u\), where \(L\) is its length and
consecutive vertices are joined by the edges of \(C\). Put
\((\bvar{u_0}\vee \bvar{u_1})\) on the edge \(u_0u_1\), and
\((\neg \bvar{u_i}\vee \bvar{u_{i+1}})\) on the edge \(u_iu_{i+1}\) for
\(1\le i\le L-1\). For \(L=1\) this is the single clause \((\bvar{u}\vee \bvar{u})\). If
\(\bvar{u}=0\), the first clause gives \(\bvar{u_1}=1\), and the others give
\(\bvar{u_2}=\dots=\bvar{u_L}=1\), which contradicts \(\bvar{u_L}=\bvar{u}=0\). Replacing
\(\bvar{u}\) by \(\neg \bvar{u}\) in these clauses gives a pattern that forces
\(\bvar{u}=0\).
\end{proof}

\begin{definition}[Barbell]\label{def:barbell}
A \emph{barbell} is a triple \(Y=(D_1,D_2,P)\), where \(D_1,D_2\) are
edge-disjoint cycles, its \emph{ends}, and \(P\) is a path from a vertex of
\(D_1\) to a vertex of \(D_2\) that uses no edge of \(D_1\) or \(D_2\). Its edge
set is \(E(Y)=E(D_1)\cup E(D_2)\cup E(P)\). If \(D_1\) and \(D_2\) share a
vertex, \(P\) may be this single vertex.
\end{definition}

\begin{lemma}[Barbells give contradictions]\label{lem:barbell-unsat}
Every barbell supports an unsatisfiable 2-CNF. Hence, if a barbell \(Y\) is a
subgraph of \(\Gin\), then with probability at least \(4^{-|E(Y)|}\) the signs on
\(E(Y)\) make \(\rform{\Gin}\) unsatisfiable.
\end{lemma}

\begin{proof}
Let \(Y=(D_1,D_2,P)\) with \(P=u_0u_1\cdots u_k\), \(u_0\in V(D_1)\),
\(u_k\in V(D_2)\). The three edge sets are disjoint, so we may choose the signs
on them separately. By Lemma~\ref{lem:forcing}, choose signs on \(D_1\) that
force \(\bvar{u_0}=1\) and signs on \(D_2\) that force \(\bvar{u_k}=0\), and put
\((\neg \bvar{u_i}\vee \bvar{u_{i+1}})\) on the edge \(u_iu_{i+1}\). A satisfying
assignment would have \(\bvar{u_0}=1\), hence \(\bvar{u_1}=\dots=\bvar{u_k}=1\), which
contradicts \(\bvar{u_k}=0\). (For \(k=0\) the two cycles force opposite values at
the same vertex.) The cycle \(D_1\) with the path \(P\) is the letter Q: its
tail \(u_k\) is forced. The second claim holds because each edge has four sign
patterns, and adding clauses keeps a formula unsatisfiable.
\end{proof}

\begin{lemma}[From walks to barbells]\label{lem:walk}
Let \(D_1,D_2\) be edge-disjoint cycles, and let \(W\) be a walk from a vertex
of \(D_1\) to a vertex of \(D_2\) that uses no edge of \(D_1\) or \(D_2\). Then
there is a path \(P\) with \(E(P)\subseteq E(W)\) such that \((D_1,D_2,P)\) is
a barbell.
\end{lemma}

\begin{proof}
Every walk contains a path between its endpoints all of whose edges are edges
of the walk: repeatedly cut out a closed subwalk between two visits of the same
vertex.
\end{proof}

\begin{lemma}[Positive excess gives two cycles]\label{lem:excess-cycles}
Every finite connected multigraph \(\Gin\) with \(\exc{\Gin}>0\) contains two
edge-disjoint cycles.
\end{lemma}

\begin{proof}
Repeatedly delete a vertex of degree \(1\) with its edge, or replace a vertex
\(w\) of degree \(2\) without self-loop and its edges \(wy,wz\) by one edge
\(yz\) (a self-loop if \(y=z\)). Each step removes a vertex, so the process
stops. Each step keeps the graph connected, raises the excess by \(\frac12\)
(one vertex and one edge fewer), and maps edge-disjoint cycles back to
edge-disjoint cycles: replace the edge \(yz\) by the two edges \(wy,wz\). The final graph \(H\) has \(\exc{H}>0\); a vertex of degree \(0\),
or of degree \(2\) with a self-loop, would be all of \(H\) and give
\(\exc{H}<0\), so all degrees are at least \(3\). Let \(n'=|V(H)|\),
\(m'=|E(H)|>\frac32n'\), let \(D\) be a shortest cycle, of length \(g\), and
suppose that \(H-E(D)\) is a forest. Then \(m'-g\le n'-1\), so
\(g>\frac{n'}2+1\). Hence \(n'\ge2\) (for \(n'=1\) we would have \(g=1\)), so
\(g\ge3\): \(H\) is simple, and \(D\) has no chord, since a chord would close
a shorter cycle. Every vertex of \(D\) has degree at least \(3\) but only two
neighbors on \(D\), so it has a neighbor outside \(D\). From
\(g>\frac{n'}2+1\) we get \(n'-g<g-2\): the \(g\) vertices of \(D\) have their
outside neighbors among fewer than \(g\) outside vertices, so by pigeonhole one
outside vertex \(w\) has two neighbors \(a\ne b\) on \(D\). Then \(wa\), \(wb\) and the shorter arc of
\(D\) from \(a\) to \(b\) form a cycle of length at most
\(\lfloor g/2\rfloor+2<g\) if \(g\ge5\). If \(g\le4\), then
\(g+1\le n'<2g-2\) (there is at least one outside vertex), which is
impossible for \(g=3\); so \(g=4\), \(n'=5\), \(w\) is the only outside
vertex, and \(w\) is adjacent to all four vertices of \(D\), in particular to
two consecutive ones, which gives a triangle. In every case this contradicts the choice of
\(D\).

Positive excess cannot be replaced by positive cycle rank, or even by
\(\cyc{\Gin}\ge2\): in a theta graph (three internally disjoint paths between two
vertices) every cycle is the union of two of the three paths, so any two
cycles share a path. If the paths have lengths \(L_1,L_2,L_3\), the
theta graph has \(L_1+L_2+L_3\) edges and \(L_1+L_2+L_3-1\) vertices, so
\(\nu=2\) but \(\xi=\frac32-\frac12(L_1+L_2+L_3)\le0\).
\end{proof}

\begin{corollary}[Positive excess gives a barbell]\label{cor:excess-barbell}
Every finite connected multigraph \(\Gin\) with \(\exc{\Gin}>0\) contains a barbell.
\end{corollary}

\begin{proof}
Take edge-disjoint cycles \(D_1,D_2\) from Lemma~\ref{lem:excess-cycles}. If
they share a vertex, we are done. Otherwise take a shortest path from
\(V(D_1)\) to \(V(D_2)\). Its inner vertices lie on neither cycle and its
endpoints lie on different cycles, so none of its edges joins two vertices of
the same cycle; hence it uses no edge of \(D_1\) or \(D_2\), and
Lemma~\ref{lem:walk} applies.
\end{proof}

%----------------------------------------------------------------------
\subsection{The barbell lemma}
%----------------------------------------------------------------------

\begin{theorem}[Barbell lemma]\label{thm:barbell}
Let \(\Gin\) be a connected multigraph that contains pairwise edge-disjoint
cycles \(C_1,\dots,C_k\). Then \(\Gin\) contains at least
\(\lceil(k-1)/3\rceil\) pairwise edge-disjoint barbells, all of whose ends are
among \(C_1,\dots,C_k\).
\end{theorem}

\begin{proof}
\emph{The auxiliary tree.} Call an edge of \(\Gin\) \emph{free} if it lies on
none of the \(C_i\). Let \(\Gamma\) be the graph on
\(V(\Gin)\cup\{w_1,\dots,w_k\}\) whose edges are the free edges of \(\Gin\) that
are not self-loops, and one edge \(w_iv\) for every \(v\in V(C_i)\). It is
connected, since an edge \(xy\) of \(C_i\) is replaced by \(x\,w_i\,y\). Fix
a spanning tree \(\mathcal T\) of \(\Gamma\) rooted at a vertex of \(\Gin\). A
\emph{cycle node} \(w_i\) has a parent \(u_i\in V(C_i)\), and its children
are distinct vertices of \(C_i\) other than \(u_i\); a tree edge between two
vertices of \(\Gin\) is a free edge. Say that a subtree \(\mathcal T_x\)
\emph{owns} its free tree edges and the edges of all \(C_i\) with
\(w_i\in\mathcal T_x\). Disjoint subtrees own disjoint edge sets, and
\(E(C_i)\) is owned by \(\mathcal T_{w_i}\) but by no subtree of a child of
\(w_i\). The \emph{top} of \(\mathcal T_x\) is \(x\) if \(x\in V(\Gin)\), and
\(u_i\) if \(x=w_i\).

\emph{The algorithm.} An \emph{unmatched cycle at} a vertex \(v\) is a pair
\((C,W)\) of a cycle \(C\in\{C_1,\dots,C_k\}\) and a walk \(W\) from a vertex
of \(C\) to \(v\) that uses no edge of \(C\). Processing \(\mathcal T\) from
the leaves up, every subtree \(\mathcal T_x\) produces barbells and passes
upward at most one unmatched cycle at its top; all of them use only edges
owned by \(\mathcal T_x\) and are pairwise edge-disjoint. The unmatched cycle
passed up by a child \(y\) of \(x\) \emph{arrives} at \(x\): if \(x,y\in
V(\Gin)\), its walk is extended by the free tree edge \(yx\), so that it ends at
\(x\); if \(y=w_i\), it already ends at \(u_i=x\); if \(x=w_i\), it ends at
\(y\in V(C_i)\). Arrivals from different children use disjoint owned edges
(plus distinct free tree edges), so the walk of one arrival avoids the cycle
of any other; and at \(x=w_i\) no arrival uses an edge of \(C_i\). By
Lemma~\ref{lem:walk}, every walk below that joins two of the cycles and
avoids both of their edges contains a barbell.

\begin{itemize}
\item \emph{\(x\in V(\Gin)\).} Pair the arrivals arbitrarily. For a pair
  \((C,W),(C',W')\), the walk \(W\overline{W'}\) (the bar denotes reversal)
  joins \(C\) to \(C'\) and avoids both, so it gives a barbell. If one
  arrival is left over, pass it up.
\item \emph{\(x=w_i\), with \(C=C_i\), \(u=u_i\), and arrivals at distinct
  vertices \(a_1,\dots,a_r\) of \(C\) other than \(u\).} If \(r=0\), pass up
  \((C,u)\), with the walk of length \(0\). If \(r=1\), the walk of the
  arrival joins its cycle to \(C\), which gives a barbell with ends
  \(C^{(1)},C\). If \(r\ge2\), then \(C\) has the three distinct vertices
  \(u,a_1,a_2\), so it is a closed path. Number the \(a_j\) in the order in
  which they are met when going around \(C\) from \(u\), put \(a_{r+1}:=u\),
  and let \(A_j\) be the arc of \(C\) from \(a_j\) to \(a_{j+1}\); the arcs
  are pairwise edge-disjoint. For \(t\le r/2\), the walk
  \(W^{(2t-1)}A_{2t-1}\overline{W^{(2t)}}\) joins \(C^{(2t-1)}\) to
  \(C^{(2t)}\) and avoids both, since \(A_{2t-1}\subseteq C\); it gives a
  barbell. If \(r\) is odd, pass up \((C^{(r)},W^{(r)}A_r)\), which ends at
  \(u\). The cycle \(C\) is \emph{discarded}.
\end{itemize}
The objects created at \(x\) use the edges of distinct arrivals, together with
either all of \(E(C_i)\) in a single object (\(r\le1\)) or distinct arcs of
\(C_i\) (\(r\ge2\)). Hence edge-disjointness and ownership are preserved.

\emph{Counting.} Every cycle \(C_i\) is, at the end, an end of a barbell, the
unmatched cycle passed up by the root, or discarded. A discarded cycle is
charged to a barbell created at the same node, and such a barbell exists
because \(r\ge2\). Different barbells are edge-disjoint and hence have
different ends. With \(b\) barbells we get \(k\le2b+1+b\), that is,
\(b\ge\lceil(k-1)/3\rceil\).
\end{proof}

\begin{remark}[Sharpness]\label{rem:sharp}
The bound is attained. Join a vertex \(y\) by an edge \(yz_i\) to each of \(t\)
gadgets, where gadget \(i\) is a triangle \(Z_i=z_ia_ib_i\) with one more
triangle \(T_i\) attached at \(a_i\) and one \(T_i'\) attached at \(b_i\).
There are \(k=3t\) edge-disjoint triangles. Suppose that all three triangles of
gadget \(i\) are ends of pairwise edge-disjoint barbells. The barbell with end
\(Z_i\) contains all of \(E(Z_i)\), and its other end is not both \(T_i\) and
\(T_i'\); say it is not \(T_i'\). The barbell with end \(T_i'\) is then a
different barbell, so it uses no edge of \(Z_i\). But its other end is a
cycle other than \(T_i'\), and every such cycle can be reached from \(T_i'\) only
through an edge of \(Z_i\): removing \(E(Z_i)\) separates \(T_i'\) from the
rest of the graph. This is a contradiction. So at most \(2t\) triangles are
ends, and there are at most \(t=\lceil(k-1)/3\rceil\) barbells. For the reduction any
bound that is linear in \(k\) for \(k\ge2\) would do; we use
\(\lceil(k-1)/3\rceil\ge k/4\) for \(k\ge2\).
\end{remark}

%----------------------------------------------------------------------
\subsection{Step 2: many edge-disjoint cycles}
%----------------------------------------------------------------------

\begin{corollary}[Many edge-disjoint cycles]\label{cor:many-cycles}
Let \(\Gin\) be a multigraph with \(\nvert\) vertices and maximum degree at most
\(\degb\) in which every component has positive excess. If \(\Gin\) contains at
least \(\alpha \nvert\) pairwise edge-disjoint cycles, where \(\alpha>0\), then
\[
  \Prx{\rform{\Gin}\SAT}\le\exp\Bigl(-\frac\alpha8\,4^{-4\degb/\alpha}\,\nvert\Bigr).
\]
\end{corollary}

\begin{proof}
Fix the cycles. Every cycle is connected and lies in one component; let
\(k_j\) be the number of fixed cycles in the \(j\)-th component. In a
component with \(k_j\ge2\), Theorem~\ref{thm:barbell} gives
\(\lceil(k_j-1)/3\rceil\ge k_j/4\) edge-disjoint barbells. In a component with
\(k_j=1\), Corollary~\ref{cor:excess-barbell} gives one barbell, and
\(1\ge k_j/4\). Barbells in different components are edge-disjoint, so we have
\(b\ge\sum_jk_j/4\ge\alpha \nvert/4\) pairwise edge-disjoint barbells.

They are edge-disjoint, so together they use at most
\(|E(\Gin)|\le\degb \nvert/2\) edges. Hence fewer than
\((\degb \nvert/2)/(4\degb/\alpha)=\alpha \nvert/8\) of them have more than
\(4\degb/\alpha\) edges, and at least \(\alpha \nvert/4-\alpha \nvert/8=\alpha \nvert/8\)
barbells have at most \(4\degb/\alpha\) edges. By Lemma~\ref{lem:barbell-unsat}, the signs on each
of them make the formula unsatisfiable with probability at least
\(q:=4^{-4\degb/\alpha}\), and these events are independent because the
barbells are edge-disjoint. Therefore
\(\Prx{\rform{\Gin}\SAT}\le(1-q)^{\alpha \nvert/8}\le e^{-q\alpha \nvert/8}\).
\end{proof}

%----------------------------------------------------------------------
\subsection{Step 3: few cycles}
%----------------------------------------------------------------------

\begin{lemma}[Non-tree balls give disjoint cycles]\label{lem:nontree-cycles}
Let \(\Gin\) have maximum degree at most \(\degb\), let \(\rad\ge1\), and suppose
that \(N\) vertices \(v\) have a ball \(\ball{\rad}{v}\) that is not a tree. Then \(\Gin\)
contains at least \(N/\ballsz{2\rad}\) pairwise vertex-disjoint cycles, each of length
at most \(2\rad+1\).
\end{lemma}

\begin{proof}
Let \(B=\ball{\rad}{v}\) be a ball that is not a tree, and let \(\mathcal T\) be a
breadth-first search tree of \(B\) from \(v\); all its vertices have depth at
most \(\rad\), because a shortest path from \(v\) stays inside \(B\). Since \(B\)
is not a tree, it has an edge \(e=xy\) outside \(\mathcal T\), possibly a self-loop or
an edge parallel to a tree edge. The edge \(e\) together with the tree paths
from \(x\) and \(y\) to their last common ancestor is a cycle in \(B\) of
length at most \(2\rad+1\).

Choose centers greedily among the \(N\) vertices, removing every candidate at
distance at most \(2\rad\) from a chosen center. Each choice removes at most
\(\ballsz{2\rad}\) candidates, so at least \(N/\ballsz{2\rad}\) centers are chosen. Their balls
of radius \(\rad\) are pairwise vertex-disjoint, and so are the cycles in them.
\end{proof}

\begin{lemma}[Removing self-loops]\label{lem:remove-loops}
Fix \(\degb\), \(\excd>0\), \(\rad\ge1\) and \(\nontree\in(0,1]\). Let \(\Gin\) have
\(\nvert\) vertices, \(2\le\dg{\Gin}{v}\le\degb\) for all \(v\), and
\(\exc{\Gin}\ge2\excd \nvert\); suppose that \(\Gin\) has at most \(\excd \nvert\)
self-loops and at most \((2\excd\nontree/\degb)\nvert\) vertices with a non-tree
ball \(\ball{\rad}{v}\). Delete all self-loops, then repeatedly delete vertices of
degree \(0\) or \(1\), and call the result \(H\). Then \(H\) satisfies \Hyp{}
with \(\degb,\excd\), it is \loctree{\rad}{\nontree},
\(|V(H)|\ge(2\excd/\degb)\nvert\), and
\(\Prx{\rform{\Gin}\SAT}\le\Prx{\rform{H}\SAT}\).
\end{lemma}

\begin{proof}
Each deleted self-loop lowers the excess by \(1\), and the later deletions do
not lower it (Lemma~\ref{lem:leaf}), so
\(\exc{H}\ge\excd \nvert\ge\excd|V(H)|\). Degrees do not grow, and the process
stops only when all degrees are at least \(2\). From
\(\excd \nvert\le\exc{H}\le|E(H)|\le\frac{\degb}2|V(H)|\) we get
\(|V(H)|\ge(2\excd/\degb)\nvert\). Since \(H\) is a subgraph of \(\Gin\), distances
in \(H\) are at least those in \(\Gin\), so the ball of radius \(\rad\) around \(v\)
in \(H\) is a connected subgraph of \(\ball{\rad}{v}\) in \(\Gin\), and it is a tree if
\(\ball{\rad}{v}\) is. Hence at most \((2\excd\nontree/\degb)\nvert\le\nontree|V(H)|\) vertices
of \(H\) have a non-tree ball. Every operation deletes clauses or keeps
satisfiability unchanged.
\end{proof}

%----------------------------------------------------------------------
\subsection{The reduction}
%----------------------------------------------------------------------

\begin{theorem}[Reduction to difficult graphs]\label{thm:reduction}
Fix \(\excpar>0\), and let \(\degb,\excd\) be as in
Proposition~\ref{prop:cleaning}. For every integer \(\rad\ge1\) and every
\(\nontree\in(0,1]\) there is \(\beta=\beta(\excpar,\rad,\nontree)>0\) such that for
every multigraph \(\Gin\) with \(\nvert\) vertices and \(\exc{\Gin}\ge\excpar \nvert\), one
of the following holds:
\begin{enumerate}[label=(\alph*)]
\item \(\Prx{\rform{\Gin}\SAT}\le e^{-\beta \nvert}\);
\item there is a multigraph \(H\) that satisfies \Hyp{} with
  \(\degb,\excd\), is \loctree{\rad}{\nontree}, has
  \(|V(H)|\ge\beta \nvert\), and satisfies
  \(\Prx{\rform{\Gin}\SAT}\le\Prx{\rform{H}\SAT}\).
\end{enumerate}
\end{theorem}

\begin{proof}
Let \(H_1\) be the graph of Proposition~\ref{prop:cleaning}, with
\(n_1=|V(H_1)|\ge(\excpar/\degb)\nvert\), and set
\[
  \alpha:=\frac{2\excd\nontree}{\degb \ballsz{2\rad}}\le\excd .
\]

\emph{Case 1: \(H_1\) has at least \(\alpha n_1\) pairwise edge-disjoint
cycles.} By Proposition~\ref{prop:cleaning}, \(H_1\) has maximum degree at most
\(\degb\) and every component has positive excess, so
Corollary~\ref{cor:many-cycles} gives
\(\Prx{\rform{H_1}\SAT}\le e^{-\beta_1n_1}\) with
\(\beta_1=\frac\alpha8\,4^{-4\degb/\alpha}\), and (a) holds for every
\(\beta\le\beta_1\excpar/\degb\).

\emph{Case 2: otherwise.} The self-loops of \(H_1\) are pairwise
edge-disjoint cycles, so there are fewer than \(\alpha n_1\le\excd n_1\) of
them. If at least \((2\excd\nontree/\degb)n_1\) vertices had a non-tree ball
\(\ball{\rad}{v}\), Lemma~\ref{lem:nontree-cycles} would give at least \(\alpha n_1\)
vertex-disjoint, hence edge-disjoint, cycles. So
Lemma~\ref{lem:remove-loops} applies to \(H_1\), and the resulting \(H\) has
\(|V(H)|\ge(2\excd/\degb)n_1\ge(2\excd\excpar/\degb^2)\nvert\). Then
(b) holds for every \(\beta\le2\excd\excpar/\degb^2\).

Taking \(\beta:=\min\{\beta_1\excpar/\degb,\,2\excd\excpar/\degb^2\}\),
which depends only on \(\excpar,\rad,\nontree\), proves the theorem.
\end{proof}

\begin{remark}[Combination with the paper and order of the constants]\label{rem:combination}
The constants are chosen in the following order: \(\excpar\) is given;
Proposition~\ref{prop:cleaning} fixes \(\degb\) and \(\excd\);
Theorem~1.12 of the paper gives \(\radz(\degb,\excd)\) and
\(\nontreeo(\degb,\excd)\), and we may assume \(\nontreeo\le1\);
Theorem~\ref{thm:reduction} is applied with \(\rad=\radz\) and \(\nontree=\nontreeo\),
which fixes \(\alpha\) and \(\beta\). In case~(b), the graph \(H\) is
difficult. Since \(\solcomp{\rform{H}}\ge1\) whenever \(\rform{H}\)
is satisfiable, the main identity (Theorem~1.10 of the paper) and
Theorem~1.12 give
\(\Prx{\rform{H}\SAT}\le\E\,\solcomp{\rform{H}}\le
2e^{-\rate|V(H)|/2}\), which is \(e^{-\Omega(\nvert)}\) because
\(|V(H)|\ge\beta \nvert\). Together with case~(a) this proves the main theorem. There
is no circularity, because \(\radz\) and \(\nontreeo\) depend only on \(\degb\)
and \(\excd\).
\end{remark}

\end{document}
